\section{VTAB：如何把“通用”变成可重复实验}

“通用”若没有 protocol，就容易变成营销词。\term{VTAB}（Visual Task Adaptation Benchmark）把评估拆成 upstream pretraining、固定 adaptation budget 与多类 downstream tasks，使研究者能够比较一个表示在广泛任务上的迁移能力，同时暴露只对某一数据域有效的捷径。

\subsection{Upstream：先冻结问题，再比较表示}

VTAB 的第一步允许研究者用自己的数据与模型学习表示。Benchmark 不强制 upstream dataset 或 architecture，但要求后续适配在统一任务与预算下进行。这样比较的对象是“整个预训练方案产生的可迁移性”，而不只是某个 backbone 的参数数量。

\lecturefigure{slide-07-vtab-upstream.jpg}{VTAB 的 upstream 阶段接收预训练数据与模型，输出待评估的表示。}{官方视频 04:40；Zhai et al. 2019。}

\begin{importantbox}{Upstream / downstream 的定义}
\term{upstream} 指预训练阶段：选择数据 $D_u$、目标 $\mathcal{L}_u$ 与参数 $\theta$。\term{downstream} 指目标任务阶段：给定小型标注集 $D_t$，学习任务头或微调参数。简化写为
\[
\theta^{\star}=\arg\min_{\theta}\mathcal{L}_u(D_u;\theta),\qquad
\phi_t^{\star}=\arg\min_{\phi_t}\mathcal{L}_t(D_t;\theta^{\star},\phi_t).
\]
第二式是否允许更新 $\theta^{\star}$，取决于 linear probe、partial fine-tuning 或 full fine-tuning protocol。
\end{importantbox}

\subsection{Task families：不能让一个 domain 主导结论}

单个平均分会隐藏任务类型差异。VTAB 把任务划分为 natural、specialized 与 structured 等组别：自然图像接近常见预训练分布；专业图像包括遥感、医疗或专用传感；结构任务强调位置、计数、深度与关系。通用表示应在多组中稳定，而不是只在语义接近的组别领先。

\lecturefigure{slide-08-vtab-task-families.jpg}{VTAB 从广泛视觉任务空间抽样 natural、specialized 与 structured tasks。}{官方视频 06:00。}

\begin{knowledgebox}{读图：平均分之前先看分组曲线}
Natural tasks 可奖励语义与纹理迁移；specialized tasks 检验 domain shift；structured tasks 更依赖空间关系。若方法只在 natural 组大幅领先，应谨慎称为“通用”。推荐报告 macro average、各组平均、最差任务、方差与 adaptation cost，而不是只给一个总分。
\end{knowledgebox}

\subsection{Adaptation：有限数据下测试表示是否可用}

Benchmark 的最后一步是在固定下游样本预算下适配并测试。这里要比较的不是“谁能在无限微调后达到最高精度”，而是表示是否已经把有用结构编码到容易被小数据利用的坐标系中。阅读 transfer protocol 时要同时记录样本数、可训练参数、调参预算和最终测试方式，因为任何一项放宽都可能把“表示更好”变成“下游训练更强”。

\lecturefigure{slide-09-vtab-transfer-protocol.jpg}{VTAB 把 upstream representation 迁移到一组下游任务，再汇总适配性能。}{官方视频 06:44。}

\begin{knowledgebox}{读图：适配器也属于实验变量}
Linear probe 测试表示的线性可分性；full fine-tuning 测试初始化是否有利；parameter-efficient adapter 则测试低成本可塑性。不同协议回答不同问题。若 A 用 full fine-tuning、B 用 linear probe，分数差异可能主要来自适配容量，而不是 upstream representation。
\end{knowledgebox}

\begin{warningbox}{Benchmark 设计本身也有 blind spots}
VTAB 提供广度，但仍是有限任务集合。它不能完整覆盖 detection、segmentation、video、open-vocabulary reasoning、calibration 与安全性。对 benchmark 优化多年后，还会出现 meta-overfitting：研究社区间接把 test suite 变成训练目标。因此结论应写成“在该 suite 与 protocol 下更可迁移”。
\end{warningbox}

\teachervoice{讲者把 VTAB 描述成前面少样本小游戏的形式化版本。老师强调的不是“19 个任务等于所有视觉”，而是先用共同 protocol 逼迫研究者回答：表示究竟在哪些任务类型上可迁移、适配成本是多少。}

\subsection{本章小结}

VTAB 将通用表示拆成 upstream、task families 与 adaptation protocol。它比单一 ImageNet accuracy 更接近“快速学习新视觉任务”，但仍需分组报告、控制适配容量并警惕 benchmark meta-overfitting。

